\begin{center}
\tagstructbegin{tag=Figure,alttext={二乘二卷积核在四乘四图像上滑动产生三乘三特征图}}
\begin{tikzpicture}[>=Latex,x=9mm,y=9mm]
\foreach \rr/\cc/\val in {
0/0/1,0/1/1,0/2/0,0/3/0,
1/0/1,1/1/1,1/2/0,1/3/0,
2/0/0,2/1/0,2/2/1,2/3/1,
3/0/0,3/1/0,3/2/1,3/3/1}{
 \filldraw[fill=paper,draw=linegray] (\cc,{3-\rr}) rectangle ++(1,1);
 \node[font=\small] at ({\cc+.5},{3.5-\rr}) {$\val$};
}
\draw[very thick,orange,rounded corners=1pt] (0,2) rectangle (2,4);
\draw[very thick,teal,dashed,rounded corners=1pt] (1,2) rectangle (3,4);
\node[font=\footnotesize,anchor=west] at (0,-.52) {橙框：首个窗口；青色虚框：右移一格};
\draw[->,blue,very thick] (4.25,2.1)--(5.22,2.1);
\node[align=center,font=\small] at (6.9,2.1)
 {$F=\begin{bmatrix}1&0\\0&-1\end{bmatrix}$\\逐格乘加};
\draw[->,blue,very thick] (8.3,2.1)--(9.1,2.1);
\node[align=center,font=\small] at (11,2.15)
 {$\begin{bmatrix}0&1&0\\1&0&-1\\0&-1&0\end{bmatrix}$\\$3\times3$ 特征图};
\end{tikzpicture}
\tagstructend

{\footnotesize 首格结果为 $1-1=0$；右移一格后为 $1-0=1$。
其余位置使用同一 $2\times2$ 核，步幅 1、不补零。}
\end{center}
\begin{center}\captionof{figure}{一个 $2\times2$ 卷积核在 $4\times4$ 图片上滑动，生成 $3\times3$ 特征图}\label{fig:cnn-window}\end{center}
